\section{Conclusiones}
\label{sec:conclusiones}

[El apartado de conclusiones debe sintetizar los hallazgos tecnicos del estudio
en una estructura numerada o tabulada. Cada conclusion debe ser una afirmacion
derivable directamente de los resultados y el analisis presentados. Evitar
introducir nueva informacion no discutida previamente.]

\begin{table}[htbp]
    \centering
    \caption{[Leyenda descriptiva de las conclusiones tecnicas.]}
    \label{tab:conclusiones}
    \setcounter{itemcount}{0}
    \begin{tabularx}{\textwidth}{@{}NX@{}}
        \toprule
        \multicolumn{1}{c}{\textbf{Item}} & \textbf{Conclusion} \\
        \midrule
        & [Conclusion 1: afirmacion tecnica derivada de resultados.] \\
        & [Conclusion 2: afirmacion tecnica derivada de resultados.] \\
        & [Conclusion 3: afirmacion tecnica derivada de resultados.] \\
        & [Conclusion 4: afirmacion tecnica derivada de resultados.] \\
        \bottomrule
    \end{tabularx}
\end{table}

[Parrafo de cierre que contextualice las conclusiones en terminos de viabilidad
tecnica, riesgos identificados y necesidad de acciones correctivas o de seguimiento.]
